% ===========================================================================
\section{The shape of a matrix}\label{sec:geometry}
\warmth{0}\quad\whisper{Every linear map is a rotation, a stretch along axes, and another rotation.}

\begin{strand}
Forget the entries for a moment and watch what $A\colon\R^n\to\R^m$ \emph{does}. Feed it the
unit sphere; out comes an \keyword{ellipsoid}. The singular value decomposition is just the
honest description of that ellipsoid: the directions of its axes, and how long they are.
Reading $A=U\Sig V\T$ right to left, $A$ is a \keyword{rotation} $V\T$, then an axis-aligned
\keyword{stretch} $\Sig$, then another rotation $U$. Three simple motions, in that order, are
\emph{all} any linear map ever does.
\end{strand}

\block{The decomposition}

\begin{definition}[singular value decomposition]\label{def:svd}
A \keyword{singular value decomposition} of $A\in\R^{m\times n}$ is a factorization
$A=U\Sig V\T$ where
\begin{itemize}[leftmargin=2em,itemsep=1pt]
  \item $U\in\R^{m\times m}$ and $V\in\R^{n\times n}$ are \keyword{orthogonal}:
        $U\T U=\fillin[1.5cm]$ and $V\T V=\fillin[1.5cm]$;
  \item $\Sig\in\R^{m\times n}$ is ``diagonal'' with entries
        $\sigma_1\ge\sigma_2\ge\cdots\ge\sigma_r>0=\sigma_{r+1}=\cdots$, the \keyword{singular values}.
\end{itemize}
The columns $u_i$ of $U$ are the \keyword{left} singular vectors; the columns $v_i$ of $V$ the
\keyword{right} singular vectors.
\end{definition}

\recall{Orthogonal matrices are exactly the rigid motions fixing the origin: they preserve every length and angle. $\det=\pm1$ (rotation or reflection).}

\block{The singular triple}
Columns of $A=U\Sig V\T$ say it most vividly. Each triple $(\sigma_i,u_i,v_i)$ satisfies
\[
  A\,v_i \;=\; \sigma_i\,u_i \qquad\text{and}\qquad A\T u_i \;=\; \sigma_i\,v_i .
\]
So $A$ takes the orthonormal directions $v_i$ to the orthogonal directions $\sigma_i u_i$:
the $v_i$ are the \emph{axes of the input circle}, the $\sigma_i u_i$ the \emph{axes of the
output ellipse}.

\begin{yourturn}
A maps the unit sphere to an ellipsoid whose semi-axis lengths are the singular values.
\begin{itemize}[leftmargin=1.3em,itemsep=1pt]
  \item For $A=\diag(3,2)$ the unit circle becomes an ellipse with semi-axes $\fillin[1cm]$ and
        $\fillin[1cm]$, along the $\fillin[2.5cm]$ directions.
  \item If some $\sigma_i=0$, the ellipsoid is \emph{flat} in that direction. Sketch what the
        unit circle maps to when $A=\diag(2,0)$. \recall{A flat ellipse = a line segment. That collapse is exactly rank loss; see §\ref{sec:subspaces}.}
\end{itemize}
\workspace[2]
\end{yourturn}

\begin{strand}
Why insist on the ordering $\sigma_1\ge\sigma_2\ge\cdots$? Because then the \emph{first} singular
triple is the single most important thing $A$ does: $\sigma_1$ is how much $A$ can stretch any
unit vector, achieved in direction $v_1$. Every later triple is the best that remains after the
earlier ones are used up. This ``greedy from the top'' structure is the whole reason truncating
the SVD gives the best simple picture (§\ref{sec:eckart}).
\end{strand}

\loose{Complex case: for $A\in\C^{m\times n}$, replace ``orthogonal/transpose'' by ``unitary/conjugate-transpose'' and every word here still holds. Singular values stay real and $\ge0$.}
